% ============================================================
%  tikz_fig_three_approaches.tex
%  三面等価（Three Equivalent Measures of GDP）
%  upLaTeX + dvipdfmx でコンパイル
% ============================================================
\documentclass[tikz, dvipdfmx]{standalone}
\usepackage{tikz}
\usetikzlibrary{arrows.meta}
\definecolor{gdpblue}  {RGB}{30,  90, 180}
\definecolor{trendred} {RGB}{190,  40,  40}
\definecolor{srasgreen}{RGB}{30,  140,  60}
\definecolor{goldcol}  {RGB}{190, 140,  10}

\begin{document}

\tikzset{
  prodbox/.style={
    draw=trendred!60, fill=trendred!7, rounded corners=6pt,
    minimum width=4.8cm, minimum height=3.0cm,
    align=center, font=\small,
  },
  expbox/.style={
    draw=gdpblue!60, fill=gdpblue!7, rounded corners=6pt,
    minimum width=4.8cm, minimum height=3.0cm,
    align=center, font=\small,
  },
  distbox/.style={
    draw=srasgreen!60, fill=srasgreen!7, rounded corners=6pt,
    minimum width=4.8cm, minimum height=3.0cm,
    align=center, font=\small,
  },
  gdpbox/.style={
    draw=goldcol!70, fill=goldcol!10, rounded corners=5pt,
    inner sep=8pt, align=center,
  },
  eqlabel/.style={
    fill=white, inner sep=4pt, font=\large\bfseries, black!55,
  },
}

\begin{tikzpicture}[>=Stealth, thick, every node/.style={font=\small}]

% ── タイトル ─────────────────────────────────────────────────
\node[font=\normalsize\bfseries, black!70] at (5.5, 9.5)
  {Three Equivalent Measures of GDP \, （三面等価）};
\draw[black!20, thin] (0.5, 9.0) -- (10.5, 9.0);

% ── 上: 生産面（Production Approach）────────────────────────
\node[prodbox] (prod) at (5.5, 7.2) {%
  \textbf{\textcolor{trendred}{Production Approach}}\\
  {\scriptsize\textcolor{trendred}{（生産面）}}\\[5pt]
  付加価値の合計\\[3pt]
  $\mathrm{GDP} = {\displaystyle\sum_i} (P_i - M_i)$\\[2pt]
  {\scriptsize\textcolor{black!45}{$P_i$: 産出額,\enspace $M_i$: 中間投入}}%
};

% ── 左下: 支出面（Expenditure Approach）─────────────────────
\node[expbox] (exp) at (0.8, 0) {%
  \textbf{\textcolor{gdpblue}{Expenditure Approach}}\\
  {\scriptsize\textcolor{gdpblue}{（支出面）}}\\[5pt]
  $Y = C + I + G + NX$\\[3pt]
  {\scriptsize\textcolor{black!45}{$C$: 消費,\enspace $I$: 投資,}}\\
  {\scriptsize\textcolor{black!45}{$G$: 政府支出,\enspace $NX$: 純輸出}}%
};

% ── 右下: 分配面（Distribution Approach）───────────────────
\node[distbox] (dist) at (10.2, 0) {%
  \textbf{\textcolor{srasgreen}{Distribution Approach}}\\
  {\scriptsize\textcolor{srasgreen}{（分配面）}}\\[5pt]
  要素所得の合計\\[3pt]
  {\scriptsize\textcolor{black!45}{賃金 $+$ 利潤}}\\
  {\scriptsize\textcolor{black!45}{$+$ 地代・利子}}\\
  {\scriptsize\textcolor{black!45}{$+$ 減価償却 $+$ 間接税}}%
};

% ── 三角形の辺 ────────────────────────────────────────────────
\draw[black!18, semithick, shorten <=12pt, shorten >=12pt]
  (prod.south west) -- (exp.north east);
\draw[black!18, semithick, shorten <=12pt, shorten >=12pt]
  (prod.south east) -- (dist.north west);
\draw[black!18, semithick, shorten <=12pt, shorten >=12pt]
  (exp.east) -- (dist.west);

% ── 辺の中点に「=」 ──────────────────────────────────────────
\node[eqlabel] at (3.1, 3.6) {$=$};
\node[eqlabel] at (7.9, 3.6) {$=$};
\node[eqlabel] at (5.5, 0.0) {$=$};

% ── 中央: GDP ボックス ────────────────────────────────────────
\node[gdpbox] at (5.5, 3.5) {%
  \textcolor{goldcol}{\Large\bfseries GDP}\\[2pt]
  {\small\textcolor{goldcol!80!black}{Gross Domestic Product}}\\
  {\scriptsize\textcolor{black!40}{（国内総生産）}}%
};

% ── 下部メッセージ ───────────────────────────────────────────
\node[font=\scriptsize, black!50, align=center] at (5.5, -1.5)
  {同じ経済活動を \textbf{3つの視点} から測定
   $\;\Rightarrow\;$ 必ず同じ値になる};

\end{tikzpicture}
\end{document}
